\section*{Discussion}

The results share a single mechanism. A Gaussian information functional is a difference of log-determinants, the log-determinant of a sample covariance matrix is a logarithmic linear spectral statistic, and the logarithm is concave. Each sample spectrum is spread around its population counterpart by the Marchenko--Pastur or Wachter law, so every log-determinant is biased downward, and by an amount that grows with dimension. Mutual information is a signed combination of three log-determinants whose dimensions sum to zero. The leading terms cancel, but the concavity does not, and what remains is the spurious information. Kullback--Leibler costs combine a log-determinant with a trace of an inverse and inherit the divergence of the latter at the left edge of the spectrum. The value of information is mutual information in other units. Random matrix theory thus provides not only the size of these effects but also their structure: exact digamma telescopes at finite $n$, Wachter integrals in the proportional limit, and phase transitions where low-rank structure meets the bulk.

Several qualifications apply. The exact results assume independent Gaussian observations, and financial returns are neither independent nor Gaussian. The placebo evidence in the application suggests that the null law is remarkably robust for block designs of moderate size. Robustness should not be assumed for heavier tails or larger dimensions, where the top of the spectrum can be dominated by a few extreme observations \citep{biroli2007top}. Elliptical and heavy-tailed versions of the theory, built on robust scatter estimators or on the general Silverstein equation with random scale, are natural extensions. The value-of-information identity is partial-equilibrium and relies on CARA preferences and joint normality. Under other preferences, mutual information remains an upper bound on the value of information for some classes of decision problems \citep{cabrales2013entropy}, but the equality is lost. The information frontier rests on the eigenvalue part of the canonical-correlation transition, which is established, and on a conjectured eigenvector counterpart, which this paper supports only by analogy with the BBP case. The dynamic analysis uses the simplest rotation-invariant model. Empirical correlation dynamics have a market mode that breaks rotation invariance, and the sufficiency result should then be applied to the bulk after projecting out the market direction.

The methodological implications are nonetheless direct. Wherever a high-dimensional Gaussian information measure is estimated, whether entropy, mutual information, transfer entropy or a Kullback--Leibler divergence between fitted Gaussian models, the exact finite-sample bias is available and should be removed. The correction is exact, needs no knowledge of the true covariance, and remains valid in the regime where the bias is largest relative to the signal. It also differs from covariance cleaning, which is the right tool for quadratic losses but biases information estimates towards zero. For empirical work on connectedness, spillovers and information transmission in finance, the corrected statistics and the exact null law give a principled way to separate genuine information from the mechanical contribution of dimension. For the economics of information, the identity between mutual information and certainty-equivalent value means that data that look valuable in high dimensions may owe much of their apparent value to dimensionality. The spurious component can be computed exactly, and it grows whenever uninformative series are added.

\section*{Conclusion}

This paper has developed the random matrix theory of Gaussian information functionals and taken it to financial data. The plug-in mutual information between two blocks of variables carries an exact spurious component that depends only on the dimensions and the sample size. It has an exact log-Beta law under independence, grows linearly in the sample size at the rate $\kappa(c_1,c_2)$ in the proportional regime, and increases strictly with every uninformative coordinate. Its spectral representation yields a closed form for the logarithmic moment of the Wachter law, and the canonical-correlation phase transition gives a frontier below which information cannot be located. An exact CARA--Gaussian identity turns all of these statements into statements about the certainty-equivalent value of information. For correlation dynamics, the Dyson eigenvalue process of a matrix Ornstein--Uhlenbeck market is informationally sufficient for its dissipation rate, which can be learned at a rate quadratic in the number of assets. On thirty-six years of US industry returns, the tools recover the market and sector modes with graded reliability, quantify the out-of-sample cost of noise dressing, and isolate an information linkage between carbon-intensive and financial industries that survives exact correction and conditioning on the market. The manuscript, data, reproduction code, simulation outputs and Lean proofs are available in the accompanying repository.
